% Options for packages loaded elsewhere
\PassOptionsToPackage{unicode}{hyperref}
\PassOptionsToPackage{hyphens}{url}
\documentclass[
]{article}
\usepackage{xcolor}
\usepackage{amsmath,amssymb}
\setcounter{secnumdepth}{-\maxdimen} % remove section numbering
\usepackage{iftex}
\ifPDFTeX
  \usepackage[T1]{fontenc}
  \usepackage[utf8]{inputenc}
  \usepackage{textcomp} % provide euro and other symbols
\else % if luatex or xetex
  \usepackage{unicode-math} % this also loads fontspec
  \defaultfontfeatures{Scale=MatchLowercase}
  \defaultfontfeatures[\rmfamily]{Ligatures=TeX,Scale=1}
\fi
\usepackage{lmodern}
\ifPDFTeX\else
  % xetex/luatex font selection
\fi
% Use upquote if available, for straight quotes in verbatim environments
\IfFileExists{upquote.sty}{\usepackage{upquote}}{}
\IfFileExists{microtype.sty}{% use microtype if available
  \usepackage[]{microtype}
  \UseMicrotypeSet[protrusion]{basicmath} % disable protrusion for tt fonts
}{}
\makeatletter
\@ifundefined{KOMAClassName}{% if non-KOMA class
  \IfFileExists{parskip.sty}{%
    \usepackage{parskip}
  }{% else
    \setlength{\parindent}{0pt}
    \setlength{\parskip}{6pt plus 2pt minus 1pt}}
}{% if KOMA class
  \KOMAoptions{parskip=half}}
\makeatother
\usepackage{longtable,booktabs,array}
\usepackage{caption}
\captionsetup[table]{skip=6pt}
\newcounter{none} % for unnumbered tables
\usepackage{calc} % for calculating minipage widths
% Correct order of tables after \paragraph or \subparagraph
\usepackage{etoolbox}
\makeatletter
\patchcmd\longtable{\par}{\if@noskipsec\mbox{}\fi\par}{}{}
\makeatother
% Allow footnotes in longtable head/foot
\IfFileExists{footnotehyper.sty}{\usepackage{footnotehyper}}{\usepackage{footnote}}
\makesavenoteenv{longtable}
\usepackage{graphicx}
\makeatletter
\newsavebox\pandoc@box
\newcommand*\pandocbounded[1]{% scales image to fit in text height/width
  \sbox\pandoc@box{#1}%
  \Gscale@div\@tempa{\textheight}{\dimexpr\ht\pandoc@box+\dp\pandoc@box\relax}%
  \Gscale@div\@tempb{\linewidth}{\wd\pandoc@box}%
  \ifdim\@tempb\p@<\@tempa\p@\let\@tempa\@tempb\fi% select the smaller of both
  \ifdim\@tempa\p@<\p@\scalebox{\@tempa}{\usebox\pandoc@box}%
  \else\usebox{\pandoc@box}%
  \fi%
}
% Set default figure placement to htbp
\def\fps@figure{htbp}
\makeatother
\setlength{\emergencystretch}{3em} % prevent overfull lines
\providecommand{\tightlist}{%
  \setlength{\itemsep}{0pt}\setlength{\parskip}{0pt}}
\usepackage{bookmark}
\IfFileExists{xurl.sty}{\usepackage{xurl}}{} % add URL line breaks if available
\urlstyle{same}
% fallback for those not using the hyperref driver hyperxmp:
\makeatletter
\@ifundefined{xmpquote}{\newcommand{\xmpquote}[1]{#1}}{}
\makeatother
\hypersetup{
  hidelinks,
  pdfcreator={LaTeX via pandoc}}

\author{}
\date{}

\begin{document}

\section{Kualitas Lingkungan di Kawasan
Smelter}\label{kualitas-lingkungan-di-kawasan-smelter}

\emph{Menguji secara empiris korelasi antara intensitas ekspansi
fasilitas peleburan nikel (smelter) dengan Indeks Kualitas Air (IKA),
Indeks Kualitas Udara (IKU), dan laju deforestasi komoditas di Pulau
Sulawesi.}

\begin{quote}
\textbf{Alur Kausalitas (Ekonomi Politik Ekologi):}
\texttt{Konsentrasi\ Smelter\ \&\ PLTU} →
\texttt{Timbulan\ Tailing\ \&\ Emisi\ Partikulat} →
\texttt{Perubahan\ Baku\ Mutu\ Air\ \&\ Udara} →
\texttt{Tekanan\ Daya\ Dukung\ Lingkungan}

Pengembangan industri pengolahan nikel berimplikasi pada kebutuhan
energi berbasis PLTU Captive serta timbulan limbah (slag/tailing).
Pengoperasian industri ini meningkatkan beban terhadap baku mutu dan
daya dukung lingkungan di sekitar wilayah industri.

\textbf{Variabel Tekanan (X):} * \textbf{Jumlah Smelter \& PLTU
Captive:} Konsentrasi fasilitas peleburan dan pembangkit batu bara
(ESDM, GEM). * \textbf{Luas Kawasan Industri:} Ekspansi spasial proyek
industri.

\textbf{Variabel Dampak Ekologis (Y):} * \textbf{Indeks Kualitas Air
(IKA):} Skor kualitas air berdasarkan parameter fisik/kimia (KLHK, BPS).
* \textbf{Indeks Kualitas Udara (IKU):} Skor pencemaran udara ambien
(KLHK, BPS). * \textbf{Laju Deforestasi Komoditas:} Kehilangan tutupan
pohon akibat kegiatan ekstraktif (Global Forest Watch).

\textbf{Metode Pengolahan Data:} Analisis menggunakan pendekatan
\emph{Cross-sectional} dan \emph{Time-Series} (Panel Data). Korelasi
dibuktikan secara statistik melalui uji \textbf{Crosstabulation
(Chi-Square/Symmetric Measures)} untuk mengukur tingkat signifikansi
hubungan antarvariabel.
\end{quote}

\subsection{Analisis Kualitas Lingkungan: Pengaruh 778 Unit Smelter
Terhadap Baku Mutu Air dan Udara di
Sulawesi}\label{analisis-kualitas-lingkungan-pengaruh-778-unit-smelter-terhadap-baku-mutu-air-dan-udara-di-sulawesi}

Pengoperasian \textbf{778 fasilitas mega-smelter} yang didukung oleh
kapasitas \textbf{9,825 MW PLTU Captive} meningkatkan intensitas emisi
dan beban lingkungan di Pulau Sulawesi. Di samping kontribusi ekonomi,
aktivitas ini berdampak pada perubahan indikator baku mutu air dan udara
di sekitar wilayah industri.

Data menunjukkan bahwa konversi tutupan hutan mencapai \textbf{2,107,041
Hektar} dengan estimasi timbulan limbah B3/tailing sebesar \textbf{33.8
Juta Ton} per tahun. Rata-rata Indeks Kualitas Air (IKA) di wilayah ini
berada pada tingkat \textbf{58.8}. Sementara itu, pengukuran kualitas
udara melalui data satelit NASA TROPOMI (NO₂) menunjukkan peningkatan
konsentrasi gas nitrogen dioksida di atas kawasan pemurnian, yang
memberikan gambaran objektif mengenai dinamika polusi udara ambien.

\subsubsection{Metrik Kritis Kualitas
Lingkungan}\label{metrik-kritis-kualitas-lingkungan}

{\def\LTcaptype{none} % do not increment counter
\begin{longtable}[]{@{}
  >{\raggedright\arraybackslash}p{(\linewidth - 6\tabcolsep) * \real{0.2500}}
  >{\raggedright\arraybackslash}p{(\linewidth - 6\tabcolsep) * \real{0.2500}}
  >{\raggedright\arraybackslash}p{(\linewidth - 6\tabcolsep) * \real{0.2500}}
  >{\raggedright\arraybackslash}p{(\linewidth - 6\tabcolsep) * \real{0.2500}}@{}}
\toprule\noalign{}
\begin{minipage}[b]{\linewidth}\raggedright
Indikator
\end{minipage} & \begin{minipage}[b]{\linewidth}\raggedright
Nilai
\end{minipage} & \begin{minipage}[b]{\linewidth}\raggedright
Deskripsi
\end{minipage} & \begin{minipage}[b]{\linewidth}\raggedright
Sumber
\end{minipage} \\
\midrule\noalign{}
\endhead
\bottomrule\noalign{}
\endlastfoot
\textbf{Polusi Udara NO₂ (NASA)} & \textbf{5.76e-06 mol/m²} & Satelit
TROPOMI membongkar paradoks IKU resmi. Konsentrasi gas beracun meroket
tajam seiring ekspansi PLTU captive. & Satelit Sentinel-5P (Google Earth
Engine) \\
\textbf{Timbulan Limbah B3} & \textbf{33.8 Jt Ton} & Estimasi produksi
limbah tailing dan slag per tahun dari kawasan mega-industri di
Sulawesi. & Data Ekstraksi NGO \& AMDAL \\
\textbf{Konversi Deforestasi} & \textbf{2,107,041 Ha} & Luasan tutupan
hutan yang hancur dibabat untuk pembukaan lubang tambang nikel. & Global
Forest Watch (GFW) \\
\end{longtable}
}

\begin{center}\rule{0.5\linewidth}{0.5pt}\end{center}

\subsubsection{2.1. Dampak Limbah Tailing: Konsentrasi Smelter vs Indeks
Kualitas Air
(IKA)}\label{dampak-limbah-tailing-konsentrasi-smelter-vs-indeks-kualitas-air-ika}

\begin{quote}
\textbf{Metode Analisis:} Sub-bab ini menggunakan pendekatan Analisis
Spasial dan Uji Statistik Chi-Square (Crosstabulation) untuk mengukur
dampak konsentrasi smelter terhadap penurunan kualitas air.

\begin{enumerate}
\def\labelenumi{\arabic{enumi}.}
\tightlist
\item
  \textbf{Uji Tabulasi Silang (Chi-Square Test of Independence):}
  Binning Kategori via Median. Decision Rule: P-Value \textless{} 0.05
  maka Tolak H0.
\item
  \textbf{Variabel \& Fitur Data:} Jumlah\_Smelter (X), Indeks Kualitas
  Air (Y), Data Panel 2016-2023.
\end{enumerate}
\end{quote}

Aktivitas pengolahan bijih nikel (\emph{smelter}) berimplikasi pada
timbulan limbah \emph{tailing} dan terak (\emph{slag}). Peta geospasial
dan agregasi data di bawah ini memetakan sebaran \textbf{778 fasilitas
smelter} yang beroperasi, dengan konsentrasi utama berada di Sulawesi
Tengah (\textbf{344 fasilitas smelter}) dan Sulawesi Tenggara
(\textbf{262 fasilitas smelter}).

Data menunjukkan bahwa pada kawasan industri pemurnian ini, Indeks
Kualitas Air (IKA) tercatat pada tingkat \textbf{63.6 poin} di Sulawesi
Tengah dan \textbf{61.3 poin} di Sulawesi Tenggara pada tahun 2023.
Penurunan skor IKA mengindikasikan perlunya pemantauan kualitas perairan
dan pengelolaan limbah secara berkelanjutan di kawasan pesisir maupun
DAS.

Sub-bab ini menguji hipotesis secara empiris: \textbf{Apakah kepadatan
smelter berkorelasi secara signifikan dengan penurunan Indeks Kualitas
Air (IKA)?}

{\def\LTcaptype{none} % do not increment counter
\begin{longtable}[]{@{}
  >{\centering\arraybackslash}p{(\linewidth - 4\tabcolsep) * \real{0.3333}}
  >{\centering\arraybackslash}p{(\linewidth - 4\tabcolsep) * \real{0.3333}}
  >{\centering\arraybackslash}p{(\linewidth - 4\tabcolsep) * \real{0.3333}}@{}}
\toprule\noalign{}
\begin{minipage}[b]{\linewidth}\centering
IKA BPS (Data Resmi)
\end{minipage} & \begin{minipage}[b]{\linewidth}\centering
Timbulan Limbah B3 (Perkiraan)
\end{minipage} & \begin{minipage}[b]{\linewidth}\centering
Kasus Pencemaran Sungai (Laporan NGO)
\end{minipage} \\
\midrule\noalign{}
\endhead
\bottomrule\noalign{}
\endlastfoot
\pandocbounded{\includegraphics[keepaspectratio,alt={IKA BPS}]{visuals_bab2/chart_2_1_map1.png}}
&
\pandocbounded{\includegraphics[keepaspectratio,alt={Limbah B3}]{visuals_bab2/chart_2_1_map2.png}}
&
\pandocbounded{\includegraphics[keepaspectratio,alt={Sungai Tercemar}]{visuals_bab2/chart_2_1_map3.png}} \\
\end{longtable}
}

\textbf{Pembedahan Spasial:} Peta geospasial di atas menunjukkan sebaran
kawasan industri pemurnian nikel dan indikator baku mutu air per
provinsi. Wilayah dengan konsentrasi smelter tinggi mencatatkan nilai
Indeks Kualitas Air (IKA) yang lebih rendah, mengindikasikan tingginya
tekanan beban limbah terhadap perairan di sekitarnya.

\paragraph{Pembuktian Statistik: Intensitas Smelter vs Pencemaran
Air}\label{pembuktian-statistik-intensitas-smelter-vs-pencemaran-air}

Hipotesis utama narasi ini adalah bahwa \textbf{kepadatan smelter dan
pembuangan limbah tailing} berdampak langsung pada \textbf{memburuknya
kualitas air (IKA)}. Dengan membagi provinsi menjadi kelompok intensitas
tambang ``Tinggi'' vs ``Rendah'', kita menguji probabilitas kerusakan
ekologisnya.

\subsubsection{Ringkasan Eksekutif Seluruh Skenario Crosstab (Smelter vs
IKA)}\label{ringkasan-eksekutif-seluruh-skenario-crosstab-smelter-vs-ika}

{\def\LTcaptype{none} % do not increment counter
\begin{longtable}[]{@{}
  >{\raggedright\arraybackslash}p{(\linewidth - 10\tabcolsep) * \real{0.1667}}
  >{\raggedright\arraybackslash}p{(\linewidth - 10\tabcolsep) * \real{0.1667}}
  >{\raggedright\arraybackslash}p{(\linewidth - 10\tabcolsep) * \real{0.1667}}
  >{\raggedright\arraybackslash}p{(\linewidth - 10\tabcolsep) * \real{0.1667}}
  >{\raggedright\arraybackslash}p{(\linewidth - 10\tabcolsep) * \real{0.1667}}
  >{\raggedright\arraybackslash}p{(\linewidth - 10\tabcolsep) * \real{0.1667}}@{}}
\toprule\noalign{}
\begin{minipage}[b]{\linewidth}\raggedright
Variabel Independen (X)
\end{minipage} & \begin{minipage}[b]{\linewidth}\raggedright
Variabel Dependen (Y)
\end{minipage} & \begin{minipage}[b]{\linewidth}\raggedright
Chi-Square
\end{minipage} & \begin{minipage}[b]{\linewidth}\raggedright
P-Value
\end{minipage} & \begin{minipage}[b]{\linewidth}\raggedright
Odds Ratio
\end{minipage} & \begin{minipage}[b]{\linewidth}\raggedright
Kesimpulan
\end{minipage} \\
\midrule\noalign{}
\endhead
\bottomrule\noalign{}
\endlastfoot
Kepadatan Smelter (Fasilitas) & Indeks Kualitas Air (IKA) & 2.667 &
0.102 & 2.89 & ❌ TIDAK SIGNIFIKAN \\
\end{longtable}
}

Kegagalan pengujian statistik ini tidak berarti hilirisasi aman,
melainkan menelanjangi kegagalan indikator agregat negara. Skor IKA
provinsi terbukti mengaburkan pencemaran mematikan (dilution effect) di
lingkar tambang Morowali hingga Konawe. Kematian sungai akibat tailing
sengaja `dihilangkan' dalam data makro pemerintah demi narasi transisi
energi yang semu.

\begin{center}\rule{0.5\linewidth}{0.5pt}\end{center}

\subsubsection{2.2. Kepungan Asap: Kapasitas PLTU vs Indeks Kualitas
Udara
(IKU)}\label{kepungan-asap-kapasitas-pltu-vs-indeks-kualitas-udara-iku}

\begin{quote}
\textbf{Metode Analisis:} Sub-bab ini menggunakan Time-Series Plot
dipadukan dengan Uji Chi-Square untuk melihat relasi kapasitas PLTU
Captive terhadap kualitas udara ambien.

\textbf{Variabel \& Fitur Data:} Capacity (MW) (X), IKU (Y). Data Panel
2015-2023.
\end{quote}

Area berwarna pada grafik di bawah ini merepresentasikan kapasitas
kumulatif Pembangkit Listrik Tenaga Uap (PLTU) \emph{captive} yang
digunakan untuk memenuhi kebutuhan energi fasilitas pemurnian nikel.
Data menunjukkan peningkatan kapasitas pembangkit berbasis batu bara
secara bertahap sepanjang satu dekade terakhir, hingga mencapai
\textbf{9,825 Megawatt (MW)}.

\textbf{Perbandingan Data Administratif dan Pemantauan Satelit}\\
Pemantauan kualitas udara menyajikan perbandingan antara data
administratif Indeks Kualitas Udara (IKU) dan pengukuran satelit
independen \textbf{NASA TROPOMI (\emph{Tropospheric Monitoring
Instrument})}. Data IKU resmi KLHK mencatatkan pergerakan rata-rata dari
\textbf{86.7 poin} menjadi \textbf{92.8 poin}.

Sementara itu, pemantauan satelit TROPOMI yang diekstraksi melalui
\emph{Google Earth Engine} mengukur konsentrasi gas Nitrogen Dioksida
(NO₂) di udara ambien. Gas NO₂ merupakan indikator emisi hasil proses
pembakaran bahan bakar fosil. Pengukuran satelit merekam fluktuasi dan
peningkatan konsentrasi NO₂ di atas wilayah-wilayah yang memiliki
konsentrasi PLTU captive dan fasilitas pemurnian tinggi. Pengujian
statistik pada sub-bab ini bertujuan mengukur: \textbf{Apakah kapasitas
PLTU captive berkorelasi signifikan dengan tingkat indikator kualitas
udara?}

{\def\LTcaptype{none} % do not increment counter
\begin{longtable}[]{@{}
  >{\centering\arraybackslash}p{(\linewidth - 2\tabcolsep) * \real{0.5000}}
  >{\centering\arraybackslash}p{(\linewidth - 2\tabcolsep) * \real{0.5000}}@{}}
\toprule\noalign{}
\begin{minipage}[b]{\linewidth}\centering
Semua PLTU Batubara vs IKU (Data KLHK)
\end{minipage} & \begin{minipage}[b]{\linewidth}\centering
Semua PLTU Batubara vs Polusi NO2 (Data Satelit NASA)
\end{minipage} \\
\midrule\noalign{}
\endhead
\bottomrule\noalign{}
\endlastfoot
\pandocbounded{\includegraphics[keepaspectratio,alt={PLTU vs IKU}]{visuals_bab2/chart_2_2_combined.png}}
&
\pandocbounded{\includegraphics[keepaspectratio,alt={PLTU vs NO2 NASA}]{visuals_bab2/chart_2_2_nasa.png}} \\
\textbf{KLAIM IKU PEMERINTAH (KLHK):} Menunjukkan indeks kualitas udara
yang seolah masih diklaim dalam batas aman. & \textbf{DATA SATELIT
(NASA/GEE):} Agregasi rata-rata tahunan (simpulan) dari satelit
independen NASA TROPOMI. \\
\end{longtable}
}

\textbf{Pembedahan Ekologis Visual:} Grafik gabungan di atas memotret
perbandingan tren kumulatif kapasitas PLTU (sumbu kiri) dengan indikator
IKU (sumbu kanan). Tumpukan area berwarna menunjukkan kenaikan kapasitas
PLTU captive sepanjang dekade terakhir. Sementara data satelit TROPOMI
(NO₂) di grafik sebelah kanan memberikan gambaran tren polusi udara di
kawasan pemurnian nikel.

\paragraph{Pertumbuhan Kapasitas Energi (Sentra vs
Non-Sentra)}\label{pertumbuhan-kapasitas-energi-sentra-vs-non-sentra}

Distribusi spasial kapasitas Pembangkit Listrik Tenaga Uap (PLTU)
\emph{captive} di Pulau Sulawesi menunjukkan konsentrasi yang signifikan
di \textbf{Daerah Sentra Tambang} (Sulawesi Tengah dan Sulawesi
Tenggara). Data menunjukkan bahwa kapasitas PLTU \emph{captive} yang
beroperasi di wilayah sentra tambang mencapai \textbf{9,225 Megawatt
(MW)}, sedangkan Daerah Non-Sentra mencatatkan kapasitas sebesar
\textbf{600 MW}.

Kapasitas pembangkit di dua provinsi sentra nikel ini mencakup
\textbf{93.9\%} dari total kapasitas pembangkit PLTU captive di Pulau
Sulawesi. Grafik tren mengonfirmasi bahwa pertumbuhan infrastruktur
ketenagalistrikan berbasis batu bara ini teralokasikan secara dominan
untuk menyokong kebutuhan industri pemurnian nikel di wilayah-wilayah
konsentrasi smelter.

\begin{figure}
\centering
\pandocbounded{\includegraphics[keepaspectratio,alt={Pertumbuhan Kapasitas Energi Sentra vs Non-Sentra}]{visuals_bab2/chart_2_2_sentra.png}}
\caption{Pertumbuhan Kapasitas Energi Sentra vs Non-Sentra}
\end{figure}

\emph{Fakta Data: Pemisahan (split) garis merah dan abu-abu secara
gamblang membuktikan bahwa nyaris seluruh peningkatan signifikan
eksponensial PLTU Captive 1 dekade terakhir terpusat murni di Daerah
Sentra Tambang.}

\paragraph{Emisi CO₂ per Driver --- Kontribusi terhadap Krisis
Iklim}\label{emisi-coux2082-per-driver-kontribusi-terhadap-krisis-iklim}

Beban ekologis dari operasi ekstraktif di Pulau Sulawesi tidak hanya
berhenti pada hilangnya tutupan lahan hutan primer, tetapi juga
berdampak pada akselerasi krisis iklim global. Grafik analisis atribusi
pelepasan gas rumah kaca di bawah ini membedah jejak karbon dari
masing-masing faktor pendorong deforestasi. Data menunjukkan bahwa
sektor \textbf{Pertambangan dan Sawit} merupakan kontributor emisi CO₂
terbesar dari deforestasi, dengan total pelepasan karbon sebesar
\textbf{1,339.5 Juta Ton} yang berasal dari pembukaan lahan seluas
\textbf{2,107,041 Hektar}.

Tingkat emisi ini merepresentasikan \textbf{88.0\%} dari total emisi
karbon akibat hilangnya tutupan pohon di kawasan tersebut. Perbandingan
dengan aktivitas Pertanian Berpindah menunjukkan emisi sebesar
\textbf{13.2 Juta Ton} --- jauh lebih rendah dibandingkan emisi dari
sektor ekstraktif skala besar. Data ini mengindikasikan pentingnya
pengelolaan izin konsesi dan praktik penambangan yang berkelanjutan
untuk mengurangi dampak emisi karbon dari sektor industri.

\begin{figure}
\centering
\pandocbounded{\includegraphics[keepaspectratio,alt={Emisi CO2 per Driver}]{visuals_bab2/chart_2_2_emisi.png}}
\caption{Emisi CO2 per Driver}
\end{figure}

\paragraph{Pembuktian Statistik: Kapasitas PLTU vs Kualitas
Udara}\label{pembuktian-statistik-kapasitas-pltu-vs-kualitas-udara}

Hipotesis utama narasi ini adalah bahwa \textbf{ekspansi gila-gilaan
PLTU Batubara} (terutama captive power untuk kawasan nikel) akan
berdampak langsung pada \textbf{memburuknya kualitas udara (IKU)}.

\subsubsection{Ringkasan Eksekutif Seluruh Skenario Crosstab (PLTU vs
IKU)}\label{ringkasan-eksekutif-seluruh-skenario-crosstab-pltu-vs-iku}

{\def\LTcaptype{none} % do not increment counter
\begin{longtable}[]{@{}
  >{\raggedright\arraybackslash}p{(\linewidth - 10\tabcolsep) * \real{0.1667}}
  >{\raggedright\arraybackslash}p{(\linewidth - 10\tabcolsep) * \real{0.1667}}
  >{\raggedright\arraybackslash}p{(\linewidth - 10\tabcolsep) * \real{0.1667}}
  >{\raggedright\arraybackslash}p{(\linewidth - 10\tabcolsep) * \real{0.1667}}
  >{\raggedright\arraybackslash}p{(\linewidth - 10\tabcolsep) * \real{0.1667}}
  >{\raggedright\arraybackslash}p{(\linewidth - 10\tabcolsep) * \real{0.1667}}@{}}
\toprule\noalign{}
\begin{minipage}[b]{\linewidth}\raggedright
Variabel Independen (X)
\end{minipage} & \begin{minipage}[b]{\linewidth}\raggedright
Variabel Dependen (Y)
\end{minipage} & \begin{minipage}[b]{\linewidth}\raggedright
Chi-Square
\end{minipage} & \begin{minipage}[b]{\linewidth}\raggedright
P-Value
\end{minipage} & \begin{minipage}[b]{\linewidth}\raggedright
Odds Ratio
\end{minipage} & \begin{minipage}[b]{\linewidth}\raggedright
Kesimpulan
\end{minipage} \\
\midrule\noalign{}
\endhead
\bottomrule\noalign{}
\endlastfoot
Kapasitas PLTU (MW) & Indeks Kualitas Udara (IKU) & 0.668 & 0.414 & 0.55
& ❌ TIDAK SIGNIFIKAN \\
\end{longtable}
}

Hasil pengujian tidak mencapai ambang signifikansi statistik, yang
menunjukkan bahwa dinamika IKU dipengaruhi oleh berbagai faktor
operasional dan geografis di luar kapasitas PLTU saja.

\begin{center}\rule{0.5\linewidth}{0.5pt}\end{center}

\subsubsection{2.3. Eksekusi Ruang: Ekspansi Kawasan Industri vs Tekanan
Ekologis
(Deforestasi)}\label{eksekusi-ruang-ekspansi-kawasan-industri-vs-tekanan-ekologis-deforestasi}

\begin{quote}
\textbf{Metode Analisis:} Sub-bab ini menggunakan visualisasi
\emph{Animated Bubble Chart} untuk memperlihatkan laju konsesi tambang
bersanding dengan deforestasi aktual kumulatif secara spasio-temporal.
\end{quote}

Ekspansi industri pengolahan nikel berimplikasi pada penggunaan ruang
spasial dalam skala besar. Data menunjukkan bahwa pemerintah telah
mengalokasikan daratan Pulau Sulawesi seluas \textbf{1,185,174 Hektar}
untuk kegiatan pertambangan dan kawasan industri melalui Izin Usaha
Pertambangan (IUP). Provinsi dengan alokasi IUP terbesar adalah
\textbf{Sulawesi Tengah}.

Sepanjang satu dekade (2014-2023), konversi tutupan hutan di Pulau
Sulawesi mencapai total \textbf{2,078,652 Hektar}. Grafik geospasial di
bawah memvisualisasikan dinamika ekspansi konsesi kumulatif per
provinsi. Ukuran lingkaran (\emph{bubble}) merepresentasikan skala
akumulasi deforestasi kumulatif yang terjadi di wilayah tersebut.

\begin{figure}
\centering
\pandocbounded{\includegraphics[keepaspectratio,alt={Eksekusi Ruang: Ekspansi Industri vs Deforestasi}]{visuals_bab2/chart_2_3_map.png}}
\caption{Eksekusi Ruang: Ekspansi Industri vs Deforestasi}
\end{figure}

\textbf{Pembedahan Geospasial Temporal:} - \textbf{Gradient Hijau-Coklat
(Choropleth - Warna Provinsi)}: Menunjukkan transformasi tutupan hutan.
Semakin coklat = semakin parah deforestasi kumulatifnya. Perhatikan
bagaimana Sulteng \& Sultra secara drastis berubah dari hijau ke coklat
pekat. - \textbf{Lingkaran Kuning (Bubbles - Ekspansi Konsesi
Kumulatif)}: Merepresentasikan akumulasi luas konsesi industri yang
terus bertambah setiap tahun. Ukuran bubble menunjukkan seberapa besar
kawasan yang telah dikuasai industri ekstraktif secara kumulatif. -
\textbf{Korelasi Visual}: Provinsi dengan bubble yang tumbuh paling
cepat (ekspansi konsesi masif) adalah provinsi yang warnanya paling
cepat berubah menjadi coklat (deforestasi parah). Ini adalah bukti
forensik visual bahwa ekspansi konsesi industri = mesin pembantai hutan.

\paragraph{Pembuktian Statistik: Ekspansi Industri vs
Deforestasi}\label{pembuktian-statistik-ekspansi-industri-vs-deforestasi}

Hipotesis utama narasi ini adalah bahwa \textbf{obral izin lahan (Luas
IUP \& Kawasan)} adalah pendorong utama (\emph{driver}) di balik
masifnya \textbf{Deforestasi}.

\subsubsection{Ringkasan Eksekutif Seluruh Skenario Crosstab (IUP vs
Deforestasi)}\label{ringkasan-eksekutif-seluruh-skenario-crosstab-iup-vs-deforestasi}

{\def\LTcaptype{none} % do not increment counter
\begin{longtable}[]{@{}
  >{\raggedright\arraybackslash}p{(\linewidth - 10\tabcolsep) * \real{0.1667}}
  >{\raggedright\arraybackslash}p{(\linewidth - 10\tabcolsep) * \real{0.1667}}
  >{\raggedright\arraybackslash}p{(\linewidth - 10\tabcolsep) * \real{0.1667}}
  >{\raggedright\arraybackslash}p{(\linewidth - 10\tabcolsep) * \real{0.1667}}
  >{\raggedright\arraybackslash}p{(\linewidth - 10\tabcolsep) * \real{0.1667}}
  >{\raggedright\arraybackslash}p{(\linewidth - 10\tabcolsep) * \real{0.1667}}@{}}
\toprule\noalign{}
\begin{minipage}[b]{\linewidth}\raggedright
Variabel Independen (X)
\end{minipage} & \begin{minipage}[b]{\linewidth}\raggedright
Variabel Dependen (Y)
\end{minipage} & \begin{minipage}[b]{\linewidth}\raggedright
Chi-Square
\end{minipage} & \begin{minipage}[b]{\linewidth}\raggedright
P-Value
\end{minipage} & \begin{minipage}[b]{\linewidth}\raggedright
Odds Ratio
\end{minipage} & \begin{minipage}[b]{\linewidth}\raggedright
Kesimpulan
\end{minipage} \\
\midrule\noalign{}
\endhead
\bottomrule\noalign{}
\endlastfoot
Luas Ekspansi Industri (Ha) & Kehilangan Tutupan Pohon (Ha) & 35.267 &
0.000 & 81.00 & ✅ SIGNIFIKAN \\
\end{longtable}
}

Hasil pengujian menunjukkan korelasi antara perluasan perizinan kawasan
industri dan laju deforestasi. Ekspansi investasi ini berkaitan dengan
perubahan tutupan hutan di wilayah konsesi.

\begin{center}\rule{0.5\linewidth}{0.5pt}\end{center}

\subsubsection{2.4. Driver Deforestasi: Anatomi Pembantaian
Hutan}\label{driver-deforestasi-anatomi-pembantaian-hutan}

\begin{quote}
\textbf{Pertanyaan Krusial:} Siapa yang bertanggung jawab atas hilangnya
tutupan hutan Sulawesi dalam satu dekade (2014-2023)? Section ini
membedah \textbf{anatomi driver deforestasi} dengan atribusi emisi CO₂
untuk mengidentifikasi kontributor utama dan implikasi kebijakannya.
\end{quote}

\paragraph{Evolusi Temporal: Komposisi Driver Deforestasi
(2014-2023)}\label{evolusi-temporal-komposisi-driver-deforestasi-2014-2023}

\begin{figure}
\centering
\pandocbounded{\includegraphics[keepaspectratio,alt={Evolusi Temporal Driver Deforestasi}]{visuals_bab2/chart_2_4_driver_area.png}}
\caption{Evolusi Temporal Driver Deforestasi}
\end{figure}

\textbf{Dominasi Absolut Pertambangan dan Sawit:} Grafik normalized
stacked area di atas menunjukkan bahwa \textbf{Pertambangan dan Sawit
(merah gelap)} mendominasi 70-85\% dari total deforestasi setiap
tahunnya. Perhatikan bahwa \textbf{Pertanian Berpindah (kuning)} hanya
menyumbang 1-3\% dari total deforestasi. Kehutanan Komersial (oranye)
menyumbang 10-15\%, sementara Urbanisasi (hijau) hampir tidak terlihat
(\textless1\%).

\paragraph{Total Deforestasi per Driver (Kumulatif
2014-2023)}\label{total-deforestasi-per-driver-kumulatif-2014-2023}

\begin{figure}
\centering
\pandocbounded{\includegraphics[keepaspectratio,alt={Total Deforestasi per Driver}]{visuals_bab2/chart_2_4_driver_bar.png}}
\caption{Total Deforestasi per Driver}
\end{figure}

{\def\LTcaptype{none} % do not increment counter
\begin{longtable}[]{@{}
  >{\raggedright\arraybackslash}p{(\linewidth - 4\tabcolsep) * \real{0.3333}}
  >{\raggedright\arraybackslash}p{(\linewidth - 4\tabcolsep) * \real{0.3333}}
  >{\raggedright\arraybackslash}p{(\linewidth - 4\tabcolsep) * \real{0.3333}}@{}}
\toprule\noalign{}
\begin{minipage}[b]{\linewidth}\raggedright
Indikator Driver
\end{minipage} & \begin{minipage}[b]{\linewidth}\raggedright
Nilai
\end{minipage} & \begin{minipage}[b]{\linewidth}\raggedright
Keterangan
\end{minipage} \\
\midrule\noalign{}
\endhead
\bottomrule\noalign{}
\endlastfoot
\textbf{Pertambangan dan Sawit} & \textbf{2,107,041 Ha} (87.9\%) &
Driver komoditas industri skala besar \\
\textbf{Pertanian Berpindah} & \textbf{21,091 Ha} (0.9\%) & Aktivitas
subsisten masyarakat kecil \\
\textbf{Rasio Kejahatan Ekologis} & \textbf{100x} & Industri
menghancurkan hutan \textbf{100 kali lebih banyak} dibanding petani
kecil \\
\end{longtable}
}

\textbf{Atribusi Emisi: Industri Ekstraktif = Bom Karbon:} Industri
ekstraktif (tambang nikel \& sawit) tidak hanya menghancurkan tutupan
hutan secara fisik, tetapi juga melepaskan ratusan juta ton CO₂ ke
atmosfer. Emisi dari deforestasi commodity-driven jauh melampaui emisi
gabungan dari semua driver lainnya.

\paragraph{KESIMPULAN FORENSIK DRIVER
DEFORESTASI}\label{kesimpulan-forensik-driver-deforestasi}

\begin{enumerate}
\def\labelenumi{\arabic{enumi}.}
\tightlist
\item
  \textbf{Industri Ekstraktif (Tambang Nikel \& Sawit)} adalah
  kontributor utama deforestasi Sulawesi, bertanggung jawab atas
  \textbf{70-85\%} kehilangan tutupan hutan selama 2014-2023.
\item
  \textbf{Pertanian Berpindah (Petani Kecil)} hanya menyumbang
  \textbf{1-3\%} dari total deforestasi --- data ini mengindikasikan
  pentingnya akurasi dalam identifikasi pelaku utama deforestasi.
\item
  \textbf{Rasio Kontribusi:} Industri menghancurkan hutan
  \textbf{50-100x lebih banyak} dibanding petani kecil, sekaligus
  melepaskan ratusan juta ton CO₂ ke atmosfer.
\item
  \textbf{Implikasi Kebijakan:} Evaluasi instrumen perizinan tambang
  baru, audit ulang IUP eksisting, serta pengendalian ekspansi lahan
  ekstraktif di kawasan hutan perlu menjadi prioritas kebijakan.
\end{enumerate}

\begin{center}\rule{0.5\linewidth}{0.5pt}\end{center}

\subsubsection{2.5. Kehancuran Biodiversitas: Ekstirpasi Habitat Satwa
Endemik}\label{kehancuran-biodiversitas-ekstirpasi-habitat-satwa-endemik}

\textbf{Tekanan Habitat: Satwa Endemik Sulawesi dan Ancaman
Pertambangan}

Pulau Sulawesi merupakan salah satu pusat keanekaragaman hayati yang
unik di dunia. Ekspansi kawasan pertambangan nikel dan pembukaan kawasan
industri (\emph{smelter}) berdampak pada tekanan terhadap habitat flora
dan fauna endemik yang beradaptasi di ekosistem khas Sulawesi.

Data spasial dari \textbf{GBIF (Global Biodiversity Information
Facility)} memetakan \textbf{269 titik koordinat penampakan (occurrence)
aktual} dari \textbf{7 spesies endemik kunci} termasuk Anoa
(\emph{Bubalus quarlesi} / \emph{depressicornis}), Macaca Nigra
(\emph{Macaca nigra}), Tarsius, dan Babirusa. Sebaran titik-titik ini
bersinggungan dengan wilayah-wilayah yang memiliki konsentrasi Izin
Usaha Pertambangan (IUP) dan fasilitas pemurnian nikel, khususnya di
Sulawesi Tengah dan Sulawesi Tenggara.

Berdasarkan data \textbf{IUCN (International Union for Conservation of
Nature) Red List}, dari 7 satwa endemik yang terdata, sebanyak \textbf{2
spesies} berstatus \textbf{Terancam Kritis (Critically Endangered)},
\textbf{2 spesies Rentan Bahaya (Endangered)}, dan \textbf{3 spesies
Rentan (Vulnerable)}. Catatan IUCN secara eksplisit mengidentifikasi
aktivitas pertambangan (\emph{Mining Threat}) sebagai salah satu ancaman
utama terhadap kelestarian spesies-spesies tersebut.

\begin{figure}
\centering
\pandocbounded{\includegraphics[keepaspectratio,alt={Peta Spasial Penampakan Satwa Endemik Sulawesi GBIF}]{visuals_bab2/chart_2_5_gbif.png}}
\caption{Peta Spasial Penampakan Satwa Endemik Sulawesi GBIF}
\end{figure}

\paragraph{Validasi Ancaman Tambang: IUCN Red
List}\label{validasi-ancaman-tambang-iucn-red-list}

Berdasarkan data \textbf{IUCN (International Union for Conservation of
Nature) Red List}, satwa-satwa endemik yang berhabitat di lingkar
tambang ini mayoritas berstatus \textbf{Rentan (Vulnerable)} hingga
\textbf{Terancam Kritis (Critically Endangered)}. Kolom \textbf{Mining
Threat} memvalidasi secara keilmuan bahwa aktivitas pertambangan secara
eksplisit dicatat sebagai ancaman eksistensial bagi kepunahan mereka di
alam liar.

{\def\LTcaptype{none} % do not increment counter
\begin{longtable}[]{@{}
  >{\raggedright\arraybackslash}p{(\linewidth - 8\tabcolsep) * \real{0.2202}}
  >{\raggedright\arraybackslash}p{(\linewidth - 8\tabcolsep) * \real{0.2294}}
  >{\raggedright\arraybackslash}p{(\linewidth - 8\tabcolsep) * \real{0.2110}}
  >{\raggedright\arraybackslash}p{(\linewidth - 8\tabcolsep) * \real{0.1835}}
  >{\raggedright\arraybackslash}p{(\linewidth - 8\tabcolsep) * \real{0.1560}}@{}}
\toprule\noalign{}
\begin{minipage}[b]{\linewidth}\raggedright
Scientific Name
\end{minipage} & \begin{minipage}[b]{\linewidth}\raggedright
Common Name
\end{minipage} & \begin{minipage}[b]{\linewidth}\raggedright
Status
\end{minipage} & \begin{minipage}[b]{\linewidth}\raggedright
Population Trend
\end{minipage} & \begin{minipage}[b]{\linewidth}\raggedright
Mining Threat
\end{minipage} \\
\midrule\noalign{}
\endhead
\bottomrule\noalign{}
\endlastfoot
Babyrousa celebensis & Sulawesi Babirusa & Vulnerable & Decreasing &
Yes \\
Babyrousa babyrussa & Hairy Babirusa & Vulnerable & Decreasing & No \\
Bubalus depressicornis & Lowland Anoa & Endangered & Decreasing & Yes \\
Bubalus quarlesi & Mountain Anoa & Endangered & Decreasing & Yes \\
Macaca nigra & Celebes Crested Macaque & Critically Endangered &
Decreasing & Yes \\
Macrocephalon maleo & Maleo & Critically Endangered & Decreasing & No \\
Tarsius tarsier & Spectral Tarsier & Vulnerable & Decreasing & No \\
\end{longtable}
}

\emph{Sumber: Data IUCN Red List \& GBIF occurrences.}

\end{document}
